\section{Выводы}
В ходе лабораторной работы реализована библиотека красно-чёрного дерева и программа-словарь с бинарной сериализацией и защитой от повреждения данных при загрузке. Программа с большим запасом удовлетворяет ограничениям задачи: на наборах из 500\,000 операций процессорное время не превышает 0.73~с (при лимите 15~с), а пиковое потребление оперативной памяти составляет до 57.2~МБ (при лимите 512~МБ).

Профилирование с помощью \texttt{gprof}, \texttt{Valgrind} (\texttt{Memcheck} и \texttt{Massif}) и \texttt{/usr/bin/time -v} позволило подтвердить отсутствие утечек памяти и точечно найти узкие места без преждевременного усложнения кода. Две компактные оптимизации --- замена парных сравнений строк на \texttt{std::string::compare} и внедрение move-семантики в \texttt{insert} --- сократили число аллокаций в куче на 20.4\% и уменьшили процессорное время на I/O-интенсивном тесте на 12.9\%. Получен практический навык построения воспроизводимых бенчмарков и интерпретации профилей CPU и динамической памяти.
\pagebreak
